\subsection{The Barangay Agent: Meso-Level Aggregation and Fiscal Execution}

The \texttt{BarangayAgent} serves as the critical meso-level administrative node within the computational architecture. It bridges the gap between the micro-level behavioral shifts of the \texttt{HouseholdAgent} and the macro-level policy interventions of the \texttt{MayorAgent}. Rather than acting as a spatial entity that moves across the grid, the Barangay Agent functions as a localized environment manager for a specific subset of households. 

To systematically detail its mechanics, this section is divided into three parts: (1) State initialization and data aggregation, (2) Autonomous baseline budget execution, and (3) The cascading incentive distribution protocol.

\paragraph{Part 1: Initialization and Meso-Level State Aggregation}

Upon initialization, the simulation instantiates seven unique \texttt{BarangayAgent} objects—representing Liangan East, Esperanza, Poblacion, Binuni, Demologan, Mati, and Babalaya. Each object is dynamically loaded with distinct demographic constraints (e.g., population sizes ranging from 463 to 1,534 households) and localized Internal Revenue Allotment (IRA) budgets derived from municipal institutional audits \citep{Delena2025}.

Computationally, the Barangay Agent serves as the primary data aggregator for the Deep Reinforcement Learning (DRL) environment. During every simulation tick, the agent queries all constituent \texttt{HouseholdAgent} objects assigned to its jurisdiction. It calculates the localized compliance rate ($R_{bgy}$) in $O(N_b)$ time, where $N_b$ is the number of households in the specific barangay:
\begin{equation}
    R_{bgy} = \frac{\sum_{i=1}^{N_b} \mathbb{I}(\text{household}_i.\text{is\_compliant})}{N_b}
    \label{eq:bgy_compliance}
\end{equation}
\addequation{Aggregated Barangay Compliance Rate}

This aggregated metric is crucial; it triggers the ``Habit Shield'' for the households (as detailed in Section 3.3.1) and is subsequently passed to the centralized \texttt{BacolodModel} to form the 17-dimensional State Space ($S_t$) evaluated by the AI Mayor \citep{TianReview2024}.

\paragraph{Part 2: Autonomous Base Layer Implementation (The Three Local Levers)}
While the Mayor Agent dynamically shifts municipal funds, the Barangay Agent continuously executes a static baseline policy using its localized budget allocation profile (e.g., 34\% IEC, 33\% Enforcement, 33\% Incentives). This ensures the community retains baseline governance even if the AI Mayor ignores them:

\begin{enumerate}
    \item \textbf{Base IEC (Micro-Interventions):} Barangay agents lack the financial capacity for mass media. Instead, their daily IEC budget is converted into grassroots activities using a greedy procurement heuristic. The algorithm prioritizes \textit{Recoridas} (roving announcements, $C_{rec} = \text{\textpeso} 200$) before allocating residual funds to \textit{Purok} meetings ($C_{pur} = \text{\textpeso} 1,000$). The resulting computational intensity applies a daily fractional scalar boost to the Attitude of all local households.
    \begin{equation}
        N_{\text{rec}} = \min\left(30, \left\lfloor \frac{Budget_{\text{IEC}}}{200} \right\rfloor \right)
        \label{eq:n_rec}
    \end{equation}
    \addequation{Barangay Recorida Frequency Calculation}

    \begin{equation}
        N_{\text{pur}} = \min\left(3, \left\lfloor \frac{\max(0, Budget_{\text{IEC}} - (N_{\text{rec}} \times 200))}{1000} \right\rfloor \right)
        \label{eq:n_pur}
    \end{equation}
    \addequation{Barangay Purok Meeting Frequency Calculation}

    \begin{equation}
        I_{\text{IEC}} = \min\left(1.0, \ 0.5\left(\frac{N_{\text{rec}}}{30}\right) + 0.5\left(\frac{N_{\text{pur}}}{3}\right)\right)
        \label{eq:i_iec}
    \end{equation}
    \addequation{Combined Local IEC Intensity Scalar}

    \noindent \textit{where:}
    \begin{itemize}
        \item $N_{\text{rec}}$: The number of \textit{Recoridas} funded and scheduled by the barangay for the given quarter.
        \item $N_{\text{pur}}$: The number of \textit{Purok} meetings funded and scheduled by the barangay.
        \item $Budget_{\text{IEC}}$: The localized financial budget strictly allocated for Information, Education, and Communication campaigns.
        \item $I_{\text{IEC}}$: The Base IEC Intensity scalar $[0.0, 1.0]$.
        \item $C_{\text{rec}}$ (200) and $C_{\text{pur}}$ (1000): The financial costs for fuel and logistical setup, respectively.
    \end{itemize}

    \item \textbf{Base Enforcement (Tanod Spawning):} To execute local enforcement, the Barangay calculates its maximum affordable headcount based on its local enforcement fund. It dynamically instantiates persistent \texttt{EnforcementAgent} objects (flagged as \texttt{is\_municipal = False}). These Tanods have a strict apprehension capacity of 1 household per patrol and suffer from a 90\% stochastic distraction rate. Their presence computes an ``Ambient Fear'' variable ($Base = 0.30$) to maintain a baseline perception of risk \citep{Zhao2022}. Cost of Enforcements is deducted from the barangays fund:
    \begin{equation}
        C_{\mathrm{Enf}} = (N_{\mathrm{Enforcers}} \times W_{\mathrm{Daily}} \times 30)
        \label{eq:cost_of_enforcement}
    \end{equation}
    \addequation{Total Local Enforcement Personnel Cost}

    \item \textbf{Base Incentives (Tier 1 Ledger):} The local incentive fund serves as the primary financial clearinghouse for positive reinforcement \citep{Chen2023}. The baseline computational cost of incentives ($C_{\mathrm{Inc}}$) is determined by the volume of successful claims:
    \begin{equation}
        C_{\mathrm{Inc}} = (V_{\mathrm{Reward}} \times N_{\mathrm{Compliant}})
        \label{eq:total_incentive_cost}
    \end{equation}
    \addequation{Total Local Incentive Distribution Cost}
\end{enumerate}

\paragraph{Part 3: The Cascading Incentive Distribution Protocol}

A major architectural feature of the \texttt{BarangayAgent} is its role as a financial clearinghouse. When a compliant \texttt{HouseholdAgent} attempts to claim a monetary reward, it invokes the barangay's \texttt{redeem\_incentive(amount)} method. To accurately simulate the ``stacking effect'' of polycentric governance \citep{Ostrom2010}, this method utilizes a \textit{cascading resource allocation tree}:

\begin{itemize}
    \item \textbf{Tier 1 (Local Funds):} The algorithm first attempts to deduct the \textpeso 500 reward from the \texttt{local\_incentive\_fund}. 
    \item \textbf{Tier 2 (Municipal Augmentation):} If local funds are exhausted, the method queries the overarching \texttt{MayorAgent}. If supplemental funds ($A_{\mathrm{Inc}, B_i}$) exist, the reward is deducted from the municipal pool.
    \item \textbf{Tier 3 (Budget Exhaustion):} If both tiers are empty, the claim is rejected, triggering a ``Perceived Unfairness'' penalty that degrades household motivation \citep{Chen2023}.
\end{itemize}

This tiered financial logic natively simulates the ``victim of success'' dynamic \citep{WorldBank2022}. If local compliance surges rapidly, the fixed municipal allocation is diluted among a larger pool of claimants, forcing the DRL agent to constantly adjust its policy to maintain systemic motivation.

\vspace{0.5cm} 
\begin{algorithm}[H]
\caption{The Barangay's Administrative and Fiscal Loop}
\label{alg:barangay_execution_english}
\begin{algorithmic}[1]
\State \textbf{Starting Information Needed:} 
\State \parbox[t]{0.9\linewidth}{$\rightarrow$ The Barangay's assigned demographic profile, local baseline budget, and a registry of all resident households.\strut}
\Statex
\State \textbf{Phase 1: Daily Aggregation and Micro-Interventions}
\State \parbox[t]{0.9\linewidth}{$\rightarrow$ \textit{Data Collection:} Query all local households to calculate the official Barangay Segregation Rate. Broadcast this rate to the AI Mayor.\strut}
\State \parbox[t]{0.9\linewidth}{$\rightarrow$ \textit{Grassroots IEC:} Spend the daily local education budget. Buy cheap \textit{Recoridas} first; if money remains, buy \textit{Purok} meetings. Apply the resulting motivation boost to all local citizens.\strut}
\State \parbox[t]{0.9\linewidth}{$\rightarrow$ \textit{Personnel Management:} Count the active Barangay Tanods on the grid. If the local budget allows for more personnel, spawn new Tanods and dispatch them into the neighborhood.\strut}
\Statex
\State \textbf{Phase 2: Processing Citizen Rewards (The Cascading Ledger)}
\If{a citizen actively segregates and requests a municipal reward}
    \If{the Barangay has its own local incentive funds remaining}
        \State \parbox[t]{0.85\linewidth}{$\rightarrow$ \textit{Approval:} Pay the citizen using local funds. The citizen remains highly motivated.\strut}
    \ElsIf{the Mayor has provided supplemental municipal funds to this Barangay}
        \State \parbox[t]{0.85\linewidth}{$\rightarrow$ \textit{Approval:} Pay the citizen using the Mayor's augmentation budget.\strut}
    \Else
        \State \parbox[t]{0.85\linewidth}{$\rightarrow$ \textit{Rejection (Budget Exhausted):} Deny the reward. The citizen feels cheated, triggering a permanent drop in their future willingness to segregate.\strut}
    \EndIf
\EndIf
\end{algorithmic}
\end{algorithm}